% !TEX program = xelatex
\documentclass{英语试卷}
\begin{document}
\examtitle{宁波市 2024 学年第二学期高考与选考模拟考试}{英语}
\par 宁波市 2024 学年第二学期高考与选考模拟考试\par
\par 英语学科 试题卷\par
\begin{notice}
\par 1. 答卷前，考生务必将自己的姓名、准考证号填写在答题卡上。\par
\par 2. 回答选择题时，选出每小题答案后，用铅笔把答题卡上对应题目的答案标号涂黑。如需改动，用橡皮擦干净后，再选涂其他答案标号。回答非选择题时，将答案写在答题卡上，写在本试卷上无效。\par
\par 3. 考试结束后，将本试卷和答题卡一并交回。\par
\end{notice}
\parttitle{第一部分 听力（共两节，满分 30 分）}
\par 做题时，先将答案标在试卷上，录音内容结束后，你将有两分钟的时间将试卷上的答案转涂到答题卡上。\par
\sectitle{第一节（共5小题：每小题1.5分，满分7.5分）}
\par 听下面 5 段对话。每段对话后有一个小题，从题中所给的 A、B、C 三个选项中选出最佳选项。听完每段对话后，你都有 10 秒钟的时间来回答有关小题和阅读下一小题。每段对话仅读一遍。\par
\question{1. What is the probable relationship between the speakers?}
\begin{choices}(3)
  \item Co-workers.
  \item Husband and wife.
  \item Fitness coach and customer.
\end{choices}
\question{2. What time is it now?}
\begin{choices}(3)
  \item 10:00 a. m.
  \item 10:10 a. m.
  \item 11:00 a. m.
\end{choices}
\question{3. Why is the man smiling?}
\begin{choices}(3)
  \item He has got a dinner invitation.
  \item He has won a large sum of money.
  \item He has been admitted to a university.
\end{choices}
\question{4. Why does the woman go to the bank?}
\begin{choices}(3)
  \item To take a survey.
  \item To open an account.
  \item To make a complaint.
\end{choices}
\question{5. What will the man do next?}
\begin{choices}(3)
  \item Sing a song.
  \item Write music.
  \item Climb a mountain.
\end{choices}
\sectitle{第二节（共15小题：每小题1.5分，满分22.5分）}
\par 听下面 5 段对话或独白。每段对话或独白后有几个小题，从题中所给的 A、B、C 三个选项中选出最佳选项。听每段对话或独白前，你将有时间阅读各个小题，每小题 5 秒钟；听完后，各小题将给出 5 秒钟的作答时间。每段对话或独白读两遍。\par
\question{6. What is the woman's opinion about online learning?}
\begin{choices}(3)
  \item She finds it just so-so.
  \item She enjoys its benefits.
  \item She questions its value.
\end{choices}
\question{7. What concerns the man most about online courses?}
\begin{choices}(3)
  \item The variety.
  \item The price.
  \item The convenience.
\end{choices}
\par 听第 7 段材料，回答第 8 至 10 题，\par
\question{8. What is the weather like?}
\begin{choices}(3)
  \item Sunny and rainy.
  \item Windy and rainy.
  \item Sunny and windy.
\end{choices}
\question{9. What are the speakers doing?}
\begin{choices}(3)
  \item Cycling.
  \item Sitting.
  \item Walking.
\end{choices}
\question{10. Who is the woman?}
\begin{choices}(3)
  \item The man's wife.
  \item The man's friend.
  \item The man's neighbor.
\end{choices}
\par 听第 8 段材料，回答第 11 至 13 题，\par
\question{11. What does the man want at first?}
\begin{choices}(3)
  \item Strawberries.
  \item Oranges.
  \item Pears.
\end{choices}
\question{12. How does the man feel recently?}
\begin{choices}(3)
  \item Angry.
  \item Excited.
  \item Uncomfortable.
\end{choices}
\question{13. What is the woman's last suggestion to the man?}
\begin{choices}(3)
  \item Seeing a doctor.
  \item Drinking fruit juice.
  \item Having a good sleep.
\end{choices}
\par 听第 9 段材料，回答第 14 至 16 题。\par
\question{14. How many camera shops are still open in town?}
\begin{choices}(3)
  \item One.
  \item Two.
  \item Three.
\end{choices}
\question{15. Why does the man think some camera shops went out of business?}
\begin{choices}(3)
  \item Cameras are old-fashioned.
  \item There are fewer photographers.
  \item People buy things on the Internet.
\end{choices}
\question{16. What will the man do this weekend?}
\begin{choices}(3)
  \item Look after his friend's kids.
  \item Take pictures for a wedding.
  \item Appreciate plants and animals.
\end{choices}
\par 听第 10 段材料，回答第 17 至 20 题，\par
\question{17. What is Bradley?}
\begin{choices}(3)
  \item A film director.
  \item A costume designer.
  \item A director's assistant.
\end{choices}
\question{18. What is Bradley doing now?}
\begin{choices}(3)
  \item Teaching skills to improve acting.
  \item Giving instructions before filming.
  \item Making rules for background actors.
\end{choices}
\question{19. How should actors handle talking in the scene?}
\begin{choices}(3)
  \item Speak out normally.
  \item Mouth words silently.
  \item Whisper to each other.
\end{choices}
\question{20. What uniform will David be wearing?}
\begin{choices}(3)
  \item A waiter uniform.
  \item A police uniform.
  \item A school uniform.
\end{choices}
\parttitle{第二部分 阅读理解（共两节，满分 50 分）}
\sectitle{第一节（共15小题：每小题2.5分，满分37.5分）}
\par 阅读下列短文，从每题所给的 A、B、C 和 D 四个选项中，选出最佳答案。\par
\passagelabel{A}
\par Promote your leadership with Stanford LEAD, our flagship online program. Gain the skills and mindset to succeed in today's complex world.\par
\par What Is LEAD?\par
\par Taught by the same world-renowned Stanford Graduate School of Business (GSB) staff who teach in our MBA program, LEAD is a year-long experiential learning journey that enables you to unlock your leadership potential. This highly flexible, adjustable, and actionable program helps you build the leadership and innovation skills needed to drive change, solve real-world challenges, and achieve your personal and professional goals.\par
\par Through Stanford GSB's innovative and academically-thorough online curriculum, you'll tackle your own real-world obstacles by integrating course concepts with daily business challenges. The program provides a unique opportunity to expand your perspectives by exchanging ideas with an incredibly diverse global group of peers, while receiving personalized, practical feedback from Stanford GSB professors, course advisors, and classmates throughout your learning journey.\par
\par Who Should Attend?\par
\par Mid-or senior-level professionals committed to adopting creative strategies to advance organizations'\par
\par leadership.\par
\par Engaged leaders passionate for development skills to better team performance.\par
\par Emerging leaders focusing on up-leveling their career path and master the foundations of impactful leadership.\par
\par Application Process for Stanford LEAD\par
\par The application process requires you to complete three components: the online application, a short essay, and an online video interview.\par
\par Fees\& Payment\par
\par Your program fee includes tuition and course materials. Fees are subject to change, as are programs, dates, and professors.\par
\par Upon admission, payment details including an online payment link will be sent by email. Full payment confirms your admission. Payment is accepted by bank wire transfer or credit card.\par
\question{21. What can we learn about LEAD?}
\begin{choices}(2)
  \item It offers on-campus courses.
  \item It is a Stanford MBA program.
  \item It focuses on theoretical learning.
  \item It aims to improve leadership skills.
\end{choices}
\question{22. Who will apply for this program?}
\begin{choices}(2)
  \item New teachers.
  \item Retired professors.
  \item Graduating students.
  \item Ambitious managers.
\end{choices}
\question{23. What is stated about Fees \& Payment?}
\begin{choices}(2)
  \item Program fees might vary.
  \item Fees exclude course materials.
  \item Cash payment is acceptable.
  \item Payment is due before application.
\end{choices}
\passagelabel{B}
\par My fingers trembled as I signed the pre-medical application form. The watercolor brushes on my desk stared back at me like abandoned friends. I'm Lora, an 18-year-old who secretly drew brain diagrams on napkins while classmates drew cartoons. "Art feeds the soul but starves the body," Mom always said. Her words haunted me as my biology textbooks slowly buried my drawing pads.\par
\par Then came the school career fair that changed everything. Dr. Eleanor Whitman from Harvard asked us to list our "hidden skills". My hand shook as I wrote "observing details"—something I’d learned from painting flower petals. "That’s medical research’s most needed skill!" she exclaimed, sharing how a student turned microbiology notes into textbook illustrations. My lab partner whispered, "Remember our photosynthesis (光合作用) comic?" For the first time, my two worlds didn't feel so separate.\par
\par The real surprise came during hospital volunteering. Watching surgeries, I noticed something—the steady hands of doctors reminded me of my brushwork. When a boy struggled to describe his pain, I drew a fire-breathing dragon on his cast. "You made medicine speak," the nurse said. Slowly, my drawing pads filled with cells that looked like colored-glass art and X-rays arranged like abstract paintings.\par
\par Mr. Dawson, our career counselor, helped connect these dots. On his whiteboard, my medical knowledge and art skills overlapped in a bright yellow circle labeled "Medical Illustration". He showed me job listings I never knew existed—graphic designers for science apps, 3D modelers for anatomy (解剖) classes. "Your ‘impractical’ skill makes you stand out," he said, pointing to my dragon drawing now used in pain management workshops.\par
\par Today, my microscope and paintbrushes share the same desk. The same hands that once hesitated between them now create 3D models of viruses for vaccine education. Sometimes I teach young patients to draw their feelings instead of describing them. My story isn't about heart and reality canceling each other out—they just\par
\par multiply possibilities.\par
\question{24. What can be inferred about Lora from paragraph 1?}
\begin{choices}(2)
  \item She had given up the hobby of drawing.
  \item She was excited about studying medicine.
  \item She felt torn between art and medical study.
  \item She regretted not practicing painting earlier.
\end{choices}
\question{25. Dr. Eleanor thought of Lora's "hidden skills" as \blanknum{}.}
\begin{choices}(2)
  \item a plus
  \item a leisure
  \item a distraction
  \item a burden
\end{choices}
\question{26. How did Lora's hospital experience influence her?}
\begin{choices}(2)
  \item It helped develop her surgical skills.
  \item It directed her toward doctor training.
  \item It boosted her enthusiasm for abstract art.
  \item It convinced her of art's role in healthcare.
\end{choices}
\question{27. What message does the author want to convey?}
\begin{choices}(2)
  \item Follow the beaten track to success.
  \item A picture is worth a thousand words.
  \item Don't put all your eggs in one basket.
  \item Two strings to your bow light the way.
\end{choices}
\passagelabel{C}
\par Global sea levels continue to rise despite efforts that cities around the world are making to reduce the effects of climate change and their carbon emissions. One firm, however, is hoping to battle the effects of climate change. Driven by alarming estimates that 36.6 million people in Asian coastal cities like Mumbai could face flooding by 2070, Luca Curci Architects, based in Venice, Italy, has released Floating City—a design that would allow urban areas to rise with the waters while making use of renewable energies to reduce their own carbon emissions.\par
\par Each 25-acre unit houses 50,000 residents on interconnected platforms, with potential expansion for 200,000 and buildings are capped at 90 meters tall for stability. "Structures must be flexible like bamboo in wind to resist waves," explains architect Curci. To achieve this, seawater resistant materials are given priority to, while hollow sections in platforms help them float. Meanwhile, movable bridges are constructed to link these "water neighborhoods" to mainland transport systems.\par
\par The project team also envisions a mix of renewable energy sources extracting power from wind, water, and the sun to power the city. During extreme weather, platforms automatically separate and regroup using intelligence algorithms inspired by how fish schools avoid hunters.\par
\par Compared to building artificial islands, this method preserves 68\% more marine (海洋) life and uses surfaces that absorb rainwater to reduce flooding risks. However, challenges persist. Seawater slowly harms materials, requiring repair every four months, and emergency plans for extreme weather are still under way. "This isn’t about abandoning coasts but evolving with them," notes designer Tim Fu.\par
\par While some question long-term possibility, the design redefines urban resilience. Unlike rigid sea walls, these dynamic structures harmonize with oceanic rhythms, offering a blueprint for climate solutions. "Cities must learn to dance with the waves," Curci says. This philosophy could transform humanity's relationship with rising seas, turning threats into sustainable opportunities.\par
\question{28. What drove Luca Curci Architects to come up with a new design?}
\begin{choices}(2)
  \item Energy crisis.
  \item Potential disasters.
  \item Tourism expansion.
  \item Overcrowded cities.
\end{choices}
\question{29. Which of the following best describes Luca Curci's design?}
\begin{choices}(2)
  \item Fictional and renewable.
  \item Adaptable and innovative.
  \item Eco-friendly and immovable.
  \item Wave-resistant and conventional.
\end{choices}
\question{30. A limitation of Luca Curci's design is that \blanknum{}.}
\begin{choices}(2)
  \item it threatens marine life
  \item it lacks clean energy options
  \item it requires regular maintenance
  \item it ignores coastal water changes
\end{choices}
\question{31. Which of the following can be a suitable title for the text?}
\begin{choices}(2)
  \item Ocean Protection: Life under Threat
  \item Rising with the Tides: Floating Cities
  \item Connecting to the Future: Eco Solutions
  \item Coastal Design: Inspiration from Bamboo
\end{choices}
\passagelabel{D}
\par It is late at night. You wander empty streets, heart pounding, muscles tense, sensing danger nearby. Now imagine the same situation, but without these responses. Would you still feel afraid? Experiences like this unveil a profound truth: our minds are not just products of the brain but emerge from its constant dialogue with the body.\par
\par At the heart of this connection is the vagus nerve (神经) —the body's longest nerve. It winds its way from the brain throughout the head and trunk, issuing commands to our organs and receiving feelings from them. Much of the confusing range of functions it regulates are automatic and operate without conscious control. These complex responses engage a group of neural networks that link brain and body. The vagus nerve is, in one way of thinking, the conduit of the mind.\par
\par The primary function of the vagus nerve is to dampen the body's responses. After a fear episode, for example, the body's powerful, life-saving threat response must be ended to restore resting heartbeat, breathing, blood flow and the rest. So powerful is its influence that the vagus nerve can literally stop the beating heart. Conversely, it can also stimulate bodily responses by releasing the brakes to accelerate.\par
\par The vagus nerve is the biological basis for many new therapies (疗法) worldwide. However, some people have taken the vagus nerve's wide-ranging bodily influence as an invitation to engage in false science. In some corners of the Internet, so-called polyvagal therapy—physical or breathing exercises that some claim reset the vagus nerve—is proposed to address just about any disorder of the mind or body. There's little evidence that these popular therapies are anything but comfort.\par
\par The vagus nerve doesn't need to be a cure for everything—its worth app....ing ... on the soil. Neuroscientists compare the vagus to an orchestra conductor, harmonizing bodily systems to compose our conscious experience. Professor Linda Rinaman once said, "We're not just brains in jars. Our thoughts and feelings are rooted in flesh and blood."\par
\question{32. How does the author introduce the topic in paragraph 1?}
\begin{choices}(2)
  \item By sharing a story.
  \item By creating a scene.
  \item By asking a question.
  \item By defining a concept.
\end{choices}
\question{33. What does the underlined word "conduit" in paragraph 2 mean?}
\begin{choices}(2)
  \item Channel.
  \item Product.
  \item Mirror.
  \item Nature.
\end{choices}
\question{34. What does paragraph 4 mainly talk about?}
\begin{choices}(2)
  \item The popularity of a biological treatment.
  \item The effectiveness of breathing exercises.
  \item The misuse of unproven nerve therapies.
  \item The evidence of vagus nerve's functions.
\end{choices}
\question{35. What do Linda Rinaman's words imply?}
\begin{choices}(2)
  \item Blood circulation determines our emotions.
  \item Physical health is unrelated to mental states.
  \item Body and mind interact to form experiences.
  \item Our brain stores feelings separate from body.
\end{choices}
\sectitle{第二节（共5小题；每小题2.5分，满分12.5分）阅读下面短文，从短文后的选项中选出能填入空白处的最佳选项。选项中有两项为多余选项。The path to healthy weight combines two key elements. Exercise and eating habits form its foundation. Some believe they demand lots of effort or planning. But that's not true. In fact, the best way to work them into our lives is by making small changes that gradually become part of our routine. \blanknum{36}}
\par Regular Exercise\par
\par Teens should get 60+ minutes of daily physical activity. \blanknum{37} All that matters is that each week you get the right balance of activity, including dancing, brisk walking, or stair-climbing. Breaking workouts into 15-minute "exercise snacks" (e. g., morning yoga, lunchtime walks) throughout the day makes goals more achievable. This strategy avoids burnout while keeping energy steady.\par
\par Good Eating Habits\par
\par \blanknum{38} Actually, crash diets cause weight rebounds. A 2022 study found 80\% of teens regained weight post-diet. Therefore, teens should eat a variety of foods, and there's nothing wrong with the occasional treat. A candy bar somehow beats snapping daily.\par
\par To build lasting habits, reshape your mindset. \blanknum{39} When stressed, try brief walks rather than reaching for snacks. Track weekly progress with a checklist—small achievements boost motivation strongly.\par
\par In conclusion, creatively design routines blending exercise and nutrition in your own way. \blanknum{40} While others choose gym workouts, you might prefer playing Frisbee with your dog. Knowing what's right for you will make it a lot easier to do!\par
\begin{choices}(1)
  \item You're all different after all.
  \item Independence matters in habit-building.
  \item View food as fuel, not emotional comfort.
  \item Low-calorie diets ensure steady weight loss.
  \item Eating well doesn't mean dieting over and over.
  \item It doesn't mean doing tough exercises or hitting the gym.
  \item Here's how to make these other healthy habits just as easy.
\end{choices}
\parttitle{第三部分 语言知识运用（共两节，满分30分）}
\sectitle{第一节（共15小题；每小题1分，满分15分）}
\par 阅读下面短文，从短文后各题所给的A、B、C和D四个选项中，选出可以填入空白处的最佳选项。\par
\par While her French classmates ate salads, Li Mei took out a package of steamed rice. When she transferred to a high school in Paris, her lunch became a daily \blanknum{41}. One Monday, one boy Pierre pointed at her lunch and whispered, "That looks \blanknum{42}!" Li Mei flushed, quietly avoiding the cafeteria.\par
\par Everything \blanknum{43} when her history teacher, Mr. Dupont, \blanknum{44} a project "Food \& Culture Week". Each student was to \blanknum{45} a traditional dish and present its cultural significance. Li Mei \blanknum{46} —would anyone care about Chinese food here? Eventually, she decided on jianbing, a flavor from hometown.\par
\par On presentation day, she \blanknum{47} the materials with trembling hands. To her surprise, Pierre volunteered to help cook. As they spread the batter (面糊), he \blanknum{48}, "I made fun of your food at first, but now I see it's part of your \blanknum{49}." When the crispy jianbing was served, classmates crowded around, asking questions about Chinese New Year traditions.\par
\par The event made everyone \blanknum{50}. Students who once \blanknum{51} her now greeted her warmly. A girl named Sophie even proposed \blanknum{52} a French-Chinese cookbook together for the school fair and suggested including family dishes passed down through generations.\par
\par On the final day, \blanknum{53} filled a book titled Bites Beyond Borders. Holding it, Li Mei realized that cultural gaps weren't about walls but doors—waiting to be unlocked. A simple \blanknum{54} had taught her more about belonging than any textbook.\par
\par "Food is the \blanknum{55} foreign policy," Mr. Dupont said, biting into a red bean croissant.\par
\cz{41}{adventure}{struggle}{routine}{necessity}
\cz{42}{messy}{ordinary}{strange}{special}
\cz{43}{mattered}{worked}{rested}{changed}
\cz{44}{praised}{postponed}{concluded}{launched}
\cz{45}{purchase}{review}{prepare}{receive}
\cz{46}{joked}{hesitated}{insisted}{protested}
\cz{47}{laid out}{held back}{heated up}{gave away}
\cz{48}{argued}{denied}{admitted}{added}
\cz{49}{identity}{interest}{personality}{quality}
\cz{50}{freed}{relaxed}{amused}{connected}
\cz{51}{betrayed}{overlooked}{envied}{misled}
\cz{52}{borrowing}{authoring}{promoting}{sponsoring}
\cz{53}{stories}{photos}{recipes}{quotes}
\cz{54}{game}{question}{gift}{lunch}
\cz{55}{tastiest}{heaviest}{oldest}{fastest}
\sectitle{第二节（共10小题；每小题1.5ladnepen分，满分15分）}
\par 阅读下面短文，在空白处填入1个适当的单词或括号内单词的正确形式。\par
\par Blue-and-white porcelain (青花瓷), a product of \blanknum{56}(share) creativity, tells a story of how East and West craftsmanship harmonized. In the Yuan Dynasty, Persian merchants brought special pigments (颜料) to Jingdezhen, where porcelain traditions \blanknum{57}(perfect). Local craftsmen mixed these pigments with white clay, creating the iconic blue patterns. By painting traditional Chinese designs using foreign materials, they achieved a perfect \blanknum{58}(combine) of traditions.\par
\par Over the centuries, this technique steadily advanced. During the Ming Dynasty, porcelain became a major export. European \blanknum{59} (noble) loved its unique charm, calling it "white gold". To meet their needs, Chinese artists added European elements like rose patterns to traditional landscapes, making the porcelain more appealing \blanknum{60} westerners.\par
\par Modern artists still use these old techniques in creative ways. At \blanknum{61} London show in 2025, a young maker used 3D printers \blanknum{62} (shape) vases but hand-painted them with dragons and Greek myths side by side. Visitors said it felt like ancient China meeting modern Europe through hybrid art.\par
\par Today, this global dialogue continues in museums worldwide. A Ming-era dish that once belonged to a French King \blanknum{63} (sit) in the Louvre. Its journey from Jingdezhen to the Palace of Versailles, marked by countless hands across continents, reflects how beauty goes beyond borders. \blanknum{64} Craig Clunas noted, "The blue that fascinates the world is never \blanknum{65} (mere) a color—it is a language."\par
\parttitle{第四部分 写作（共两节，满分40分）}
\sectitle{第一节（满分15分）}
\par 假定你是李华，上周你参加了志愿者社团组织的“帮助老人跨越数字鸿沟”活动，请你写一篇短文在班级英语展示角分享此次活动，内容包括：\par
\par (1) 活动内容;\par
\par (2) 你的体会。\par
\par 注意：\par
\par (1) 写作词数应为 80 左右;\par
\par (2) 请按如下格式在答题纸的相应位置作答。\par
\par Bridging the Digital Divide for the Elderly\par
\sectitle{第二节（满分25分）}
\par 阅读下面材料，根据其内容和所给段落开头语续写两段，使之构成一篇完整的短文。\par
\par The old workshop door creaked open. Clara walked in quietly. Her shoes stepped on wood shavings (刨花) that sparkled in the sunlight. Through the dusty windows, the morning light shone on Grandfather's rough hands. He was gluing a maple chair leg.\par
\par The wall calendar still showed January. Its yellowed pages were curling, just like dry leaves. An open order book lay on the desk. February's page was empty. There had been no new orders since mid-January. Clara held her breath and recorded Grandfather with her phone. She focused on his quick hands as he measured angles with a metal ruler. "This might help him," she thought, eager to show how skilled he was.\par
\par However, when she posted the video online, negative comments flooded in. "Old ways for a dying business!" "No wonder the business is failing!" Saddened, Clara deleted the video quickly. A spider crawled over the empty order book. Its blank pages showed not many people wanted handmade things these days.\par
\par "Want to see my treasure?" Grandpa’s voice broke the silence. He opened a leather album filled with pencil sketches (草图) like library chairs from 1958 and park benches from 1962. A faded photo slipped out: a smiling boy missing two teeth was grinning at the camera, holding a curved (弯曲的) spoon proudly.\par
\par "Timmy, my former student," Grandfather smiled, touching the curve of the spoon in the photo. "Now he designs furniture for tall buildings. Last month's newspaper had an article about his curved stairs. It was inspired by this very spoon." He tapped the photo, adding, "Machines make things fast, but hands remember stories."\par
\par Clara looked at the photo—the curved spoon, Timmy's smile, and Grandpa's young handwriting on the back: Timmy's first repair, age 7. Beside her, fresh wood shavings were shining. "Stories must be told," Clara murmured, taking out her phone again. This time, she aimed to share not just the work, but the story behind it.\par
\par 注意：\par
\par (1) 续写词数应为 150 左右;\par
\par (2) 请按如下格式在答题卡的相应位置作答。\par
\writingpara{Paragraph 1: The next morning, Clara uploaded a new video. \blanknum{}}
\writinglines{8}
\writingpara{Paragraph 2: A week later, Timmy, now a famous furniture designer, walked in the workshop. \blanknum{}}
\writinglines{8}
\end{document}
